% ============================================================================
% Глава 2: Теоретична и технологична обосновка
% ============================================================================

\section{ТЕОРЕТИЧНА И ТЕХНОЛОГИЧНА ОБОСНОВКА}

\subsection{Нерелационни бази данни}

Нерелационните бази данни (NoSQL) представляват клас системи за управление на данни, които се различават от традиционните релационни бази данни по няколко ключови характеристики \cite{kleppmann2017designing}. Основните особености на NoSQL системите включват:

\begin{itemize}
    \item \textbf{Липса на фиксирана схема}: За разлика от релационните бази данни, NoSQL системите позволяват съхранение на данни без предварително дефинирана схема. Това дава възможност за лесно добавяне на нови полета и промяна на структурата на данните без необходимост от миграции.
    \item \textbf{Разпределено съхранение}: NoSQL базите са проектирани за работа в разпределени среди, като поддържат хоризонтално мащабиране чрез добавяне на нови възли към клъстера \cite{cap_theorem}.
    \item \textbf{Различни модели на данни}: NoSQL системите предлагат разнообразни модели - документо-ориентирани, ключ-стойност, колонно-ориентирани и графови бази данни.
    \item \textbf{Производителност при големи обеми}: Оптимизирани са за обработка на големи количества данни и висока скорост на четене/запис.
\end{itemize}

\subsection{Документо-ориентиран модел на данни}

За настоящия проект е избран документо-ориентираният модел на NoSQL бази данни, който е реализиран чрез MongoDB. Този модел е най-подходящ за платформата за персонализирано обучение поради следните причини:

\begin{enumerate}
    \item \textbf{Естествено представяне на йерархични данни}: Образователните данни имат естествена йерархична структура (курсове → модули → лекции, студенти → прогрес). JSON-подобният формат на документите в MongoDB естествено отразява тази сложност без необходимост от нормализация.

    \item \textbf{Гъвкавост при еволюция на схемата}: В образователните платформи често се добавят нови типове оценки, алтернативни методи за завършване или персонализирани характеристики. Документо-ориентираният модел позволява лесно разширение на структурата на данните без прекъсване на работата на системата.

    \item \textbf{Висока производителност при четене}: Аналитичните заявки и персонализираните препоръки, които са критични за платформата за обучение, се изпълняват значително по-бързо в MongoDB благодарение на възможността за съхранение на свързани данни в един документ.

    \item \textbf{Подходящ за времеви серии}: Данните за активността на студентите представляват времеви серии с различна честота на записване. MongoDB предлага оптимизации за работа с такива данни.
\end{enumerate}

\subsection{MongoDB - избрана технология}

MongoDB е водеща документо-ориентирана NoSQL база данни, която предоставя \cite{mongodb_manual}:

\begin{itemize}
    \item \textbf{Документо-ориентирано съхранение}: Данните се съхраняват като BSON (Binary JSON) документи, които могат да съдържат вложени обекти, масиви и различни типове данни.

    \item \textbf{Богат заявков език}: Поддържа сложни заявки с филтриране, сортиране, агрегиране и текстови търсения. Aggregation Framework позволява изграждане на ETL тръбопроводи за трансформация на данни.

    \item \textbf{Индексиране и производителност}: Поддържа различни типове индекси (B-tree, геопространствени, текстови, TTL) за оптимизация на заявките \cite{mongodb_indexing}.

    \item \textbf{Репликация и шардинг}: Осигурява висока наличност чрез репликация и хоризонтално мащабиране чрез шардинг.

    \item \textbf{Вградени инструменти за управление}: MongoDB Compass за визуално управление, mongosh за интерактивна работа, и MongoDB Atlas за облачно разполагане.
\end{itemize}

MongoDB е избран поради своята зрялост, активна общност и богат набор от функции, които напълно покриват изискванията на проекта.

\subsection{Използвани технологии}

Проектът е разработен с помощта на съвременни технологии, които осигуряват надеждност, производителност и лесна поддръжка:

\subsubsection{Node.js и Express.js}

За реализацията на приложния слой е избран Node.js - платформа за изпълнение на JavaScript извън браузъра \cite{nodejs_docs}. В комбинация с Express.js framework се осигурява:

\begin{itemize}
    \item \textbf{Неблокиращ I/O}: Асинхронният модел позволява ефективна обработка на множество едновременни връзки, което е критично за образователни платформи с висока активност на потребителите.

    \item \textbf{Единен език за цялата система}: JavaScript се използва както за бекенд, така и за фронтенд, което опростява разработката и поддръжката.

    \item \textbf{Богата екосистема}: NPM предоставя достъп до хиляди модули, включително официалния MongoDB драйвер за Node.js.
\end{itemize}

\subsubsection{Допълнителни технологии}

\begin{itemize}
    \item \textbf{Jest}: Framework за unit тестване, който осигурява 100\% покритие на кода и автоматизирана проверка на функционалностите.

    \item \textbf{Docker}: Контейнеризация на приложението за лесно разполагане и консистентност между различни среди \cite{docker_docs}.

    \item \textbf{GitHub Actions}: CI/CD платформа за автоматизирано тестване и деплоймънт при всяка промяна в кода.
\end{itemize}

\subsection{Обосновка на технологичния стек}

Изборът на технологичен стек е направен въз основа на специфичните изисквания на образователните платформи:

\begin{enumerate}
    \item \textbf{Съответствие с домейна}: MongoDB е широко използвана в образователни приложения поради своята гъвкавост и производителност при работа с йерархични данни.

    \item \textbf{Производителност}: Node.js осигурява висока производителност при обработка на множество едновременни потребителски сесии.

    \item \textbf{Лесна интеграция}: Всички компоненти (MongoDB, Node.js, Docker) работят добре заедно и имат добра поддръжка в облачни среди.

    \item \textbf{Бъдеща разширяемост}: Архитектурата позволява лесно добавяне на нови функционалности, микросървиси и интеграция с външни системи.
\end{enumerate}

Тази технологична основа предоставя солидна платформа за разработка на надеждни и мащабируеми образователни системи с NoSQL бази данни.
